% Options for packages loaded elsewhere
\PassOptionsToPackage{unicode}{hyperref}
\PassOptionsToPackage{hyphens}{url}
\documentclass[
]{article}
\usepackage{xcolor}
\usepackage{amsmath,amssymb}
\setcounter{secnumdepth}{-\maxdimen} % remove section numbering
\usepackage{iftex}
\ifPDFTeX
  \usepackage[T1]{fontenc}
  \usepackage[utf8]{inputenc}
  \usepackage{textcomp} % provide euro and other symbols
\else % if luatex or xetex
  \usepackage{unicode-math} % this also loads fontspec
  \defaultfontfeatures{Scale=MatchLowercase}
  \defaultfontfeatures[\rmfamily]{Ligatures=TeX,Scale=1}
\fi
\usepackage{lmodern}
\ifPDFTeX\else
  % xetex/luatex font selection
\fi
% Use upquote if available, for straight quotes in verbatim environments
\IfFileExists{upquote.sty}{\usepackage{upquote}}{}
\IfFileExists{microtype.sty}{% use microtype if available
  \usepackage[]{microtype}
  \UseMicrotypeSet[protrusion]{basicmath} % disable protrusion for tt fonts
}{}
\makeatletter
\@ifundefined{KOMAClassName}{% if non-KOMA class
  \IfFileExists{parskip.sty}{%
    \usepackage{parskip}
  }{% else
    \setlength{\parindent}{0pt}
    \setlength{\parskip}{6pt plus 2pt minus 1pt}}
}{% if KOMA class
  \KOMAoptions{parskip=half}}
\makeatother
\usepackage{color}
\usepackage{fancyvrb}
\newcommand{\VerbBar}{|}
\newcommand{\VERB}{\Verb[commandchars=\\\{\}]}
\DefineVerbatimEnvironment{Highlighting}{Verbatim}{commandchars=\\\{\}}
% Add ',fontsize=\small' for more characters per line
\newenvironment{Shaded}{}{}
\newcommand{\AlertTok}[1]{\textcolor[rgb]{1.00,0.00,0.00}{\textbf{#1}}}
\newcommand{\AnnotationTok}[1]{\textcolor[rgb]{0.38,0.63,0.69}{\textbf{\textit{#1}}}}
\newcommand{\AttributeTok}[1]{\textcolor[rgb]{0.49,0.56,0.16}{#1}}
\newcommand{\BaseNTok}[1]{\textcolor[rgb]{0.25,0.63,0.44}{#1}}
\newcommand{\BuiltInTok}[1]{\textcolor[rgb]{0.00,0.50,0.00}{#1}}
\newcommand{\CharTok}[1]{\textcolor[rgb]{0.25,0.44,0.63}{#1}}
\newcommand{\CommentTok}[1]{\textcolor[rgb]{0.38,0.63,0.69}{\textit{#1}}}
\newcommand{\CommentVarTok}[1]{\textcolor[rgb]{0.38,0.63,0.69}{\textbf{\textit{#1}}}}
\newcommand{\ConstantTok}[1]{\textcolor[rgb]{0.53,0.00,0.00}{#1}}
\newcommand{\ControlFlowTok}[1]{\textcolor[rgb]{0.00,0.44,0.13}{\textbf{#1}}}
\newcommand{\DataTypeTok}[1]{\textcolor[rgb]{0.56,0.13,0.00}{#1}}
\newcommand{\DecValTok}[1]{\textcolor[rgb]{0.25,0.63,0.44}{#1}}
\newcommand{\DocumentationTok}[1]{\textcolor[rgb]{0.73,0.13,0.13}{\textit{#1}}}
\newcommand{\ErrorTok}[1]{\textcolor[rgb]{1.00,0.00,0.00}{\textbf{#1}}}
\newcommand{\ExtensionTok}[1]{#1}
\newcommand{\FloatTok}[1]{\textcolor[rgb]{0.25,0.63,0.44}{#1}}
\newcommand{\FunctionTok}[1]{\textcolor[rgb]{0.02,0.16,0.49}{#1}}
\newcommand{\ImportTok}[1]{\textcolor[rgb]{0.00,0.50,0.00}{\textbf{#1}}}
\newcommand{\InformationTok}[1]{\textcolor[rgb]{0.38,0.63,0.69}{\textbf{\textit{#1}}}}
\newcommand{\KeywordTok}[1]{\textcolor[rgb]{0.00,0.44,0.13}{\textbf{#1}}}
\newcommand{\NormalTok}[1]{#1}
\newcommand{\OperatorTok}[1]{\textcolor[rgb]{0.40,0.40,0.40}{#1}}
\newcommand{\OtherTok}[1]{\textcolor[rgb]{0.00,0.44,0.13}{#1}}
\newcommand{\PreprocessorTok}[1]{\textcolor[rgb]{0.74,0.48,0.00}{#1}}
\newcommand{\RegionMarkerTok}[1]{#1}
\newcommand{\SpecialCharTok}[1]{\textcolor[rgb]{0.25,0.44,0.63}{#1}}
\newcommand{\SpecialStringTok}[1]{\textcolor[rgb]{0.73,0.40,0.53}{#1}}
\newcommand{\StringTok}[1]{\textcolor[rgb]{0.25,0.44,0.63}{#1}}
\newcommand{\VariableTok}[1]{\textcolor[rgb]{0.10,0.09,0.49}{#1}}
\newcommand{\VerbatimStringTok}[1]{\textcolor[rgb]{0.25,0.44,0.63}{#1}}
\newcommand{\WarningTok}[1]{\textcolor[rgb]{0.38,0.63,0.69}{\textbf{\textit{#1}}}}
\setlength{\emergencystretch}{3em} % prevent overfull lines
\providecommand{\tightlist}{%
  \setlength{\itemsep}{0pt}\setlength{\parskip}{0pt}}
\usepackage{bookmark}
\IfFileExists{xurl.sty}{\usepackage{xurl}}{} % add URL line breaks if available
\urlstyle{same}
\hypersetup{
  hidelinks,
  pdfcreator={LaTeX via pandoc}}

\author{}
\date{}

\begin{document}

\section{BOINC: A Platform for Volunteer
Computing}\label{boinc-a-platform-for-volunteer-computing}

\textit{[Image omitted: images/787f399e322a7e2bab8897e9268ff311a580715f4e4f15e34d6e3c23de9d175d.jpg]}

David P. Anderson

Received: 21 December 2018 / Accepted: 20 October 2019 / Published
online: 16 November 2019 © The Author(s) 2019

Abstract ``Volunteer computing'' is the use of consumer digital devices
for high-throughput scientific computing. It can provide large computing
capacity at low cost, but presents challenges due to device
heterogeneity, unreliability, and churn. BOINC, a widely-used
open-source middleware system for volunteer computing, addresses these
challenges. We describe BOINC's features, architecture, implementation,
and algorithms.

Keywords BOINC · Volunteer computing · Distributed computing ·
Scientific computing · High-throughput computing

\section{1 Introduction}\label{introduction}

\section{1.1 Volunteer Computing}\label{volunteer-computing}

Volunteer computing (VC) is the use of consumer digital devices, such as
desktop and laptop computers, tablets, and smartphones, for
high-throughput scientific computing. Device owners participate in VC by
installing a program that downloads and executes jobs from servers
operated by science projects. There are currently about 30 VC projects
in many scientific areas and at many institutions. The research enabled
by VC has resulted in numerous papers in Nature, Science, PNAS, Physical
Review, Proteins, PloS Biology, Bioinformatics, J. of Mol. Biol., J.
Chem. Phys, and other top journals {[}1{]}.

About 700,000 devices are actively participating in VC projects. These
devices have about 4 million CPU cores and 560,000 GPUs, and
collectively provide an average throughput of 93 PetaFLOPS. The devices
are primarily modern, high-end computers: they average 16.5 CPU
GigaFLOPS and 11.4 GB of RAM, and most have a GPU capable of
general-purpose computing using OpenCL or CUDA.

The potential capacity of VC is much larger: there are more than 2
billion consumer desktop and laptop computers {[}2{]}. Current models
have a peak performance (including GPU) of over 3 TeraFLOPS, giving a
total peak performance of 6000 ExaFLOPS. The capacity of mobile devices
is similar: there are about 10 billion mobile devices, and current
mobile processors have on-chip GPUs providing as much as 4 TeraFLOPS
{[}3{]}. Studies suggest that of people who learn about VC, between 5\%
and 10\% participate, and that desktop and mobile devices are available
to compute for VC about 60\% and 40\% of the time respectively {[}4,
5{]}. Taking these factors into account, the near-term potential
capacity of VC is on the order of hundreds of ExaFLOPS.

The monetary cost of VC is divided between volunteers and scientists.
Volunteers pay for buying and maintaining computing devices, for the
electricity to power these devices, and for Internet connectivity.
Scientists pay for a server and for the system administrators,
programmers and web developers needed to operate the VC project.
Operating a typical VC project involves

a few Linux server computers and a part-time system administrator,
costing the research project on the order of \$100 K per year. Several
BOINC projects (Einstein@Home, Rosetta@home, SETI@home) are of this
scale, and they average about 2 PetaFLOPS throughput each.

How much would this computing power cost on a commercial cloud? As of
October 2019, a compute-intensive Amazon EC2 instance type
(``c4.8xlarge'') provides 640 GFLOPS {[}6{]}, so it would take 3125 such
instances to provide 2 PetaFLOPS. Instances of this type cost \$1.59 per
hour {[}7{]}, for a total cost of \$43.52 M per year -- 435 times the
cost of using VC. The cost is lower for ``spot'' instances: \$0.30 per
hour {[}8{]}, for a total cost of \$8.21 M/year. However, the
unreliability of spot instances may cause problems for long-running
jobs. For projects with GPU applications, AWS's ``p3.2xlarge'' instance
type, with a NVIDIA Tesla V100 GPU, provides 7.104 TeraFLOPS {[}9{]} and
costs \$3.06/h, for a total cost of \$7.56 M/year. Spot instances of
this node type cost \$0.95/h, for a total cost of \$2.34 M/year. In all
cases, AWS also charges for data and storage. Thus, VC is potentially
much cheaper than EC2; the numbers are similar for other cloud
providers. Kondo et al.~compared the cost of volunteer and commercial
cloud computing in more detail, and reached similar conclusions
{[}10{]}.

In terms of power consumption, data-center nodes are typically more
efficient (i.e.~have greater FLOPS/watt) than consumer devices. However,
to compare net energy usage, two other factors must be considered.
First, data centers use air conditioning to remove the heat generated by
hardware; consumer devices generally do not. Second, in cool climates,
the heat generated by consumer devices contributes to ambient heating,
so the net energy cost of computing may be zero. Thus it's not clear
whether VC is globally more or less energy efficient than data-center
computing. In any case, this does not affect VC's cost to the scientist,
since the energy is paid for by volunteers.

VC is best suited to high-throughput computing (HTC), where workloads
consist of large groups or streams of jobs, and the goal is high rate of
job completion rather than low job turnaround time. VC is less suited to
workloads that have extreme memory or storage requirements, or for which
the ratio of network communication (i.e.~input and output file size) to
computing is extremely high.

VC differs from other forms of HTC, such as grids and clouds, in several
ways. Each of these factors presents challenges that a VC platform must
address.

\begin{itemize}
\tightlist
\item
  VC's computers are anonymous, untrusted, inaccessible, and
  uncontrollable. They may misbehave in any way, and misbehavior can't
  be stopped or punished; it must be detected and ignored. In
  particular, computers may return (intentionally or not) incorrect
  computational results, so the correctness of these must be validated
  in some way.\\
\item
  The computers are heterogeneous in all hardware and software
  dimensions {[}11{]}. This requires either multiple application
  versions or the use of virtualization. It also complicates the
  estimation of job runtimes.\\
\item
  Creating the resource pool involves recruiting and retaining
  volunteers. This requires incentive features such as teams, computing
  credit accounting and screensaver graphics.\\
\item
  The scale of VC is larger: up to millions of computers and millions of
  jobs per day. Hence the server software must be efficient and
  scalable.
\end{itemize}

\section{1.2 BOINC}\label{boinc}

Most VC projects use BOINC, an open-source middleware system for VC
{[}12{]}. BOINC lets scientists create and operate VC projects, and lets
volunteers participate in these projects. Volunteers install an
application (the BOINC client) and then choose one or more projects to
support. The client is available for desktop platforms (Windows, Mac,
Linux) and for mobile devices running Android. BOINC is designed to
compute invisibly to the volunteer. It runs jobs at the lowest process
priority and limits their total memory footprint to prevent excessive
paging. On mobile devices it runs jobs only when the device is plugged
in and fully charged, and it communicates only over WiFi.

BOINC can handle essentially all HTC applications. Many BOINC projects
run standard scientific programs such as Autodock, Gromacs, Rosetta,
LAMMPS, and BLAST. BOINC supports applications that use GPUs (using CUDA
and OpenCL), that use multiple CPUs (via OpenMP or OpenCL), and that run
in virtual machines or Docker containers.

The development of BOINC began in 2002, and was carried out by a team at
UC Berkeley led by the author,

with funding from the National Science Foundation. The first BOINC-based
project was launched in 2004. BOINC is distributed under the open-source
LGPL v3 license, but applications need not be open-source.

This paper describes BOINC and how it addresses the challenges listed
above. Section 2 describes the overall structure of BOINC. Section 3
focuses on its job-processing features, while Section 4 discusses job
failures and timeouts. Section 5 describes the client and server
software. Scheduling is discussed in Section 6. Mechanisms for
incentivizing and supporting volunteers are described in Sections 7 and
8. Section 9 presents related work, and Section 10 offers conclusions
and suggestions for future work.

\section{2 The High-Level Structure of
BOINC}\label{the-high-level-structure-of-boinc}

BOINC's architecture involves multiple components interacting through
HTTP-based RPC interfaces. It is designed to be modular and extensible.

Figure 1 shows the components of BOINC; these components are described
subsequently. The shaded boxes represent the BOINC software
distribution; unshaded boxes represent components developed by third
parties.

\section{2.1 Projects and Volunteers}\label{projects-and-volunteers}

In BOINC terminology, a project is an entity that uses BOINC to do
computing. Each project has a server, based on the BOINC server
software, that distributes jobs to volunteer devices. Projects are
autonomous and independent. A project may serve a single research group
(SETI@home), a set of research groups using a common set of applications
(Climateprediction.net, nanoHUB@home) or a set of unrelated scientists
(IBM World Community Grid, BOINC@TACC). BOINC provides HTTP-based remote
job submission APIs so that projects can submit and manage jobs using
web-based interfaces (e.g., science gateways) or via existing systems
such as HTCondor.

Projects may operate a public-facing web site. The BOINC server software
provides database-driven web scripts that support volunteer login,
message boards, profiles, leader boards, badges, and other community and
incentive functions. The project may extend this with their own web
content, such as descriptions and news of their research. A project is
identified by the URL of its web site.

A volunteer is the owner of a computing device (desktop, laptop, tablet,
or smartphone) who wants to participate in VC. They do so by a)
installing the BOINC client on the device; b) selecting projects to
support and creating an account on each project, and c) attaching the
BOINC client to each account. Each attachment can be assigned a resource
share indicating the relative amount of computing to be done, over the
long term, for the project.

\section{2.2 Client/Server Interaction}\label{clientserver-interaction}

When attached to a project, the BOINC client periodically issues RPCs to
the project's server to report completed jobs and get new jobs. The
client then downloads application and input files, executes jobs, and
uploads

\textit{[Image omitted: images/774a239309a5a65cf1ae501b72632244af451dd121c2328109377fcc494e8d8f.jpg]}

flowchart

\begin{Shaded}
\begin{Highlighting}[]
\NormalTok{graph TD}
\NormalTok{    A["Volunteer device"] {-}{-}\textgreater{}|user control| B["BOINC client"]}
\NormalTok{    B {-}{-}\textgreater{}|get, report jobs| C["BOINC project server"]}
\NormalTok{    C {-}{-}\textgreater{}|create accounts| D["account managers"]}
\NormalTok{    C {-}{-}\textgreater{}|get attached projects| B}
\NormalTok{    C {-}{-}\textgreater{}|submit, control jobs manage files| E["remote job submission"]}
\NormalTok{    C {-}{-}\textgreater{}|get host, user credit data| F["credit statistics web sites"]}
\NormalTok{    B {-}{-}\textgreater{}|get project list, client versions| C}
\end{Highlighting}
\end{Shaded}

Fig. 1 The components and RPC interfaces of the BOINC system

output files. All communication uses client-initiated HTTP; thus the
BOINC client works behind firewalls that allow outgoing HTTP, and with
HTTP proxies. The server must be publicly accessible via HTTP.

Downloaded files can be compressed in transit; the client uncompresses
them. Files need not be located on project servers. The validity of
downloaded files is checked using hashes, and (for application files)
code signing; see Section 3.8.

Files are uploaded by an HTTP POST operation, which is handled by a CGI
program on the server. Upload filenames include a random string to
prevent spoofing. To prevent DoS attacks, projects can optionally use
encryption-based tokens to limit the size of uploaded files. File
uploads need not go to the project server; for example,
Climateprediction.net collaborates with scientists in different
countries, and the result files are sent directly to the collaborators'
servers.

The client also interacts occasionally with a server operated by the
BOINC project itself, to obtain current lists of projects and of client
versions.

All client/server interactions handle failure using exponential back-off
in order to limit the rate of requests when a server resumes after a
period of being off-line.

\section{2.3 Account Managers}\label{account-managers}

The method described above for finding and attaching projects can be
cumbersome when a) the volunteer has many computers and wants to attach
or detach projects on all of them, or b) the volunteer wants to attach
many projects.

To address this, BOINC provides a framework for web-based account
managers (AMs). The BOINC client, rather than being explicitly attached
to projects, can be attached to an AM. The client periodically does an
RPC to the AM. The reply includes a list of projects (and corresponding
accounts) to which the client should attach. The client then attaches to
these projects and detaches from others.

The AM architecture was designed to support ``project selection'' web
sites, which show a list of projects with checkboxes for attaching to
them. This solves the two problems described above; for example, a
volunteer with many computers can attach them all to the AM account, and
can then change project attachments across all the computers with a
single AM interaction. Two such account managers have been developed:
GridRepublic (https://www.gridrepublic.org/) and BAM!
(https://boincstats.com/en/bam/). More recently, the AM architecture has
been used to implement the coordinated VC model (Section 10.1), in which
volunteers choose science areas rather than projects.

When a volunteer selects a project on an AM, they initially have no
account on that project. The AM needs to be able to create an account.
To support this, projects export a set of HTTP RPCs for creating and
querying accounts.

\section{2.4 Computing and Keyword
Preferences}\label{computing-and-keyword-preferences}

Volunteers can specify computing preferences that determine how their
computing resources are used. For example, they can enable CPU
throttling, in which computing is turned on and off with a configurable
duty cycle, in order to limit CPU heat. Other options include whether to
compute when the computer is in use; limits on disk and memory usage;
limits on network traffic; time-of-day limits for computing and network,
how much work to buffer, and so on. Computing preferences can be set
through a web interface on either a project or account manager; they
propagate to all the volunteer's computers that are attached to that
account. Preferences can also be set locally on a computer.

More recently, we have introduced the concept of keyword preferences.
BOINC defines two hierarchies of ``keywords'': one set for science areas
(Physics, Astrophysics, Biomedicine, cancer research, etc.) and another
set for location of the research project (Asia, United States, U.C.
Berkeley, etc.) Volunteers can mark keywords as ``yes'' or ``no'',
indicating either that they prefer to do computing tagged with the
keyword, or that they are unwilling to do it. This system has been used
for two purposes:

\begin{itemize}
\tightlist
\item
  Projects that do research in multiple science areas can let their
  volunteers choose what areas to support. The project assigns keywords
  to jobs. In selecting jobs to send to a given volunteer, BOINC will
  prefer jobs with ``yes'' keywords, and won't choose jobs with ``no''
  keywords.\\
\item
  It enables the ``coordinated VC model'' (see Section 10.1), in which
  account manager dynamically assigns projects to volunteers in a way
  that reflects project keywords and volunteer keyword preferences.
\end{itemize}

\section{3 Job Processing Abstractions and
Features}\label{job-processing-abstractions-and-features}

\section{3.1 Handling Resource
Heterogeneity}\label{handling-resource-heterogeneity}

The pool of volunteered computers is highly diverse {[}11{]}. At a
hardware level, computers may have Intel-compatible, Alpha, MIPS, or ARM
CPUs, and their 32- and 64-bit variants. Intel-compatible CPUs can have
various special instructions like SSE3. ARM processors may have NEON or
VFP floating-point units. The computers run various operating systems
(Windows, Mac OS X, Linux, FreeBSD, Android) and various versions of
these systems. There are a range of GPUs made by NVIDIA, AMD, and Intel,
with various driver versions offering different features.

BOINC provides a framework that lets projects use as many of these
computing resources as possible, and to use them as efficiently as
possible. Projects can build versions of their programs for specific
combinations of hardware and software capabilities. These are called app
versions, and the collection of them for a particular program is called
an app. Jobs are submitted to apps, not app versions.

BOINC defines a set of platforms, which are generally the combination of
a processor type and operating system: for example, (Windows, Intelx64).
Each app version is associated with a platform, and is sent only to
computers that support that platform.

For finer-grained control of version selection, BOINC provides the plan
class mechanism. A plan class is a function that takes as input a
description of a computer (hardware and software), and returns a)
whether an app version can run on that computer; b) if so what resources
(CPU and GPU, possibly fractional) it will use, and c) the peak FLOPS of
those resources. For example, a plan class could specify that an app
version requires a particular set of GPU models, with a minimum driver
version. Plan classes can be specified either in XML or as C++
functions. App versions can optionally be associated with a plan class.

App versions have integer version numbers. In general, new jobs are
processed using the latest version for a given (platform, plan class)
pair; in-progress jobs may be using older versions.

When the BOINC client requests jobs from a server, it includes a list of
platforms it supports, as well as a description of its hardware and
software. When the server dispatches a job to a computer, it selects an
app version with a matching platform, and whose plan class function
accepts the computer. The details of this scheduling policy are
described in Section 6.4.

As an example, SETI@home has 9 supported platforms, 2 apps, 86 app
versions, and 34 plan classes.

\section{3.2 Jobs Properties}\label{jobs-properties}

A BOINC job includes a reference to an app and a collection of input
files. A job has one or more instances. Jobs have input files; instances
have output files. When the server dispatches an instance to a client,
it specifies a particular app version to use. The client downloads the
input files and the files that make up the app version, executes the
main program, and (if it completes successfully) uploads the resulting
output files and the standard error output.

In the following we use ``FLOPs'' (lower-case s) as the plural of FLOP
(floating-point operation). FLOPS (upper-case s) denotes FLOPs per
second.

Each job has a number of properties supplied by the submitter. These
include:

\begin{itemize}
\tightlist
\item
  An estimate of the number of FLOPs performed by the job. This is used
  to predict runtime; see Section 6.3.\\
\item
  A maximum number of FLOPs, used to abort jobs that go into an infinite
  loop.\\
\item
  An estimate of RAM working-set size, used in server job selection (see
  Section 6.4).\\
\item
  An upper bound on disk usage, used both for server job selection and
  to abort buggy jobs that use unbounded disk space.\\
\item
  An optional set of keywords describing the job's science area and
  origin.
\end{itemize}

\section{3.3 Result Validation}\label{result-validation}

Hosts may return incorrect results because of hardware errors or
malicious user behavior. Projects typically require that the final
results of the job are correct with high probability. There may be
application-specific ways of doing this: for example, for physical
simulations one could check conservation of energy and the stability of
the final configuration. For other cases, BOINC provides a mechanism
called replication-based validation in which each job is processed on
two unrelated computers. If the outputs agree, they are accepted as
correct; otherwise a third instance is created and run.

This is repeated until either a quorum of consistent instances is
achieved, or a limit on the number of instances is reached.

What does it mean for two results to agree? Because of differences in
floating-point hardware and math libraries, two computers may return
different but equally valid results. McIntosh et al.~{[}13{]} showed
that by carefully selecting compilers, compiler options, and libraries
it's possible to get bitwise-identical results for FORTRAN programs
across the major platforms. In general, however, discrepancies will
exist. For applications that are numerically stable, these discrepancies
lead to small differences in the final results. BOINC allows projects to
supply application-specific validator functions that decide if two
results are equivalent, i.e.~whether corresponding values agree within
specified tolerances.

Physical simulations are typically not stable. For such applications,
BOINC provides a mechanism called homogeneous redundancy {[}14{]} in
which computers are divided into equivalence classes, based on their
hardware and software, such that computers in the same class compute
identical results for the app. Once a job instance has been dispatched
to a computer, further instances are dispatched only to computers in the
same equivalence class. BOINC supplies two equivalence relations: a
coarse one involving only OS and CPU vendor, and a finer-grained one
involving CPU model as well. Projects can define their own equivalence
relations if needed.

In some cases there are discrepancies between app versions; for example,
CPU and GPU versions may compute valid but incomparable results. To
address this, BOINC provides an option called homogeneous app version.
Once an instance has been dispatched using a particular app version,
further instances use only the same app version. This can be used in
combination with homogeneous redundancy.

Basic replication-based validation reduces effective computing capacity
by a factor of at least two. BOINC provides a refinement called adaptive
replication that moves this factor close to one. The idea is to identify
hosts that consistently compute correct results (typically the vast
majority of hosts) and use replication only occasionally for jobs sent
to these hosts. Actually, since some computers are reliable for CPU jobs
but unreliable for GPU jobs, we maintain this ``reputation'' at the
granularity of (host, app version).

Adaptive replication works as follow: the BOINC server maintains, for
each (host, app version) pair, a count N of the number of consecutive
jobs that were validated by replication. Once N exceeds a threshold,
jobs sent to that host with that app version are replicated only some of
the time; the probability of replication goes to zero as N increases.
Adaptive replication can achieve a low bound on the error rate
(incorrect results accepted as correct), even in the presence of
malicious volunteers, while imposing only a small throughput overhead.

Result validation is useful for preventing credit cheating (Section 7).
Credit is granted only to validated results, and ``faked'' results
generally will fail validation. For applications that usually return the
same answer (e.g.~primality checks), cheaters could simply always return
this result, and would usually get credit. To prevent this, the
application can add extra output that depends on intermediate results of
the calculations.

\section{3.4 The BOINC Runtime
Environment}\label{the-boinc-runtime-environment}

We now discuss the environment in which the BOINC client runs project
applications, and how the client controls and communicates with running
jobs.

The BOINC client installers for Windows and Mac OS X create an
unprivileged user account under which BOINC runs applications. This
provides a measure of protection against buggy or malicious
applications, since they are unable to read or write files belonging to
existing users.

The BOINC client stores its files in a dedicated directory. Within this
directory, there is a project directory for each project to which the
client is attached, containing the files associated with that project.
There are also a number of job directories, one for each in-progress
job. Job directories contain symbolic links to files in the
corresponding project directory, allowing multiple concurrent jobs to
share single copies of executable and input files.

To enforce volunteer preferences, to support CPU throttling, and to
support the client GUI's monitoring and control functions, the client
must be able to suspend, resume, and abort jobs, and to monitor their
CPU time and memory usage. Jobs may consist of a number of processes.
Some modern operating systems provide the needed process-group
operations, but not all do, and BOINC must run on older systems as well.
So BOINC implements these functions itself, based on message-passing
between the client and the application. This message-passing uses shared
memory, which is supported by all operating systems. There are queues
for

messages from client to application (suspend, resume, quit) and from
application to client (current CPU time, CPU time at last checkpoint,
fraction done, working set size). Each side polls its incoming queues
every second.

All applications run under BOINC must implement this message protocol.
One way to do this is to rebuild the application for the target
platforms, linking it with a BOINC runtime library, and using its API;
we call this a native application. Another way is to use an existing
executable together with a wrapper program, as described in the next
section. These two approaches are shown in Fig. 2.

For sequential programs, the runtime library creates a timer thread that
manages communication with the client, while the original worker thread
runs the program. On Windows, the timer thread controls the worker
thread with Windows thread primitives. The Unix thread library doesn't
let one thread suspend another, so we use a periodic signal, handled by
the worker thread, and suspend the worker by having this signal handler
sleep.

For multithread applications (e.g., OpenMP or OpenCL CPU) the process
structure is a bit different. On Unix, the program forks. The parent
process handles process control messages, which it implements using
signals. The child process runs the program, and uses a timer thread to
handle status and checkpointing messages.

BOINC supports application-level checkpoint/re-start. Every few minutes
the client sends a message asking the application to checkpoint. When
the application reaches a point where it can efficiently checkpoint
(e.g.~the start of its outer loop) it writes a checkpoint file and sends
a message to the client indicating this. Thus the client knows when
applications have checkpointed, and can avoid preempting jobs that
haven't checkpointed in a long time, or ever.

In applications that use GPUs, the CPU part of the program dispatches
short ``kernels'' to the GPU. It's imperative that the application not
be suspended while a kernel is in progress. So the BOINC runtime library
supports masked sections during which suspension and abortion are
deferred. Execution of GPU kernels (as well as writing checkpoint files)
should be done in a masked section.

By default, the BOINC client runs applications at the lowest process
priority. Non-CPU-intensive and wrapper applications are run at normal
priority. GPU applications have been found to perform poorly if run at
low priority, so they're also run at normal priority.

CPU throttling (see Section 2.4) is implemented by having the client
suspend and resume the application, with one-second granularity, and
with a duty cycle determined by the volunteer's setting.

Applications may encounter a transient situation where they can't
continue; for example, a GPU application may fail to allocate GPU
memory. The BOINC API provides a temporary exit function that causes the
application to exit and tells the client to schedule it again after a
certain amount of time. Jobs that do this too many times are aborted.

Applications should make minimal assumptions about host software. If
they need libraries other than standard system libraries, they must
include them. When the BOINC client runs a job, it modifies its library
search path to put the job directory first, followed by the project
directory.

Fig. 2 The process structure of the BOINC runtime environment\\
\textit{[Image omitted: images/56f0c738a917fd87780f336d36c8331587d4421682fe6b3a8dd970a75441e6a9.jpg]}

flowchart

\begin{Shaded}
\begin{Highlighting}[]
\NormalTok{graph TD}
\NormalTok{    A["client"] {-}{-}\textgreater{}|share{-}memory message queues| B["worker thread"]}
\NormalTok{    A {-}{-}\textgreater{}|share{-}memory message queues| C["timer thread"]}
\NormalTok{    B {-}{-}\textgreater{}|thread control| C}
\NormalTok{    C {-}{-}\textgreater{}|thread control| B}
\NormalTok{    D["wrapped application"] {-}{-}\textgreater{}|process or VM control| E["application"]}
\NormalTok{    D {-}{-}\textgreater{}|wrapped| F["wrapper"]}
\NormalTok{    F {-}{-}\textgreater{}|wrapped| E}
\end{Highlighting}
\end{Shaded}

\section{3.5 Application Graphics}\label{application-graphics}

App versions can include a second executable program that generates a
visualization of the computation in progress. The graphics application
may be run by the BOINC screensaver (which cycles among the graphics of
running jobs) or by a command from the BOINC Manager.

The graphics application for a job is run in the job's directory. It may
get state information by reading files in that directory, or by creating
a shared memory segment shared with the application. Graphics programs
typically use OpenGL for portability.

\section{3.6 Wrapped and Virtual Machine
Applications}\label{wrapped-and-virtual-machine-applications}

Native applications must be modified to use the BOINC API, and then
rebuilt for the target platforms. This may be difficult or impossible,
so BOINC provides ``wrapper'' mechanisms for using existing executables.

The BOINC wrapper allows arbitrary executables to be run under BOINC, on
the platforms for which they were built. The wrapper acts as an
intermediary between the BOINC client and the application. It handles
runtime messages from the client (Section 3.4) and translates them into
OS-specific actions (e.g.~signals on Unix) to control the application.
BOINC provides versions of the wrapper for the major platforms. An app
version includes the appropriate version of the wrapper (as the main
program), and the application executable. The wrapper can be instructed
to run a sequence of executables. Marosi et al.~developed a more general
wrapper that supports arbitrary workflows using a shell-like scripting
language {[}15{]}.

BOINC supports applications that run in virtual machines using
VirtualBox, a multi-platform open-source VM system for Intel-compatible
computers. Applications can then be developed on Linux and run on
Windows and Mac OS X computers with no additional work. The BOINC client
detects and reports the presence of VirtualBox. VM apps are only
dispatched to hosts where VirtualBox is installed. On Windows, the
recommended BOINC installer installs VirtualBox as well.

To support VM apps, BOINC provides a VBox wrapper program, which
interfaces between the BOINC client and the VirtualBox executive. It
translates runtime system messages to VirtualBox operations. VirtualBox
supports sharing of directories between host and guest systems. The VBox
wrapper creates a shared directory for VM jobs, places input files
there, and retrieves output files from there.

In addition to reducing heterogeneity problems, VM technology provides a
strong security sandbox, allowing untrusted applications to be used. It
also provides an application-independent checkpoint/restart capability,
based on the VirtualBox ``snapshot'' feature. The wrapper tells
VirtualBox to create a snapshot every few minutes. If computing is
stopped (e.g.~because the host is turned off) the job can later be
restarted, on that host, from the most recent snapshot.

BOINC's VM system can be used together with container systems like
Docker {[}16{]}. Docker-based applications consist of a Docker image
(the OS and library environment of choice) and an executable to run in
the container. The container is run within a VirtualBox virtual machine.
Docker images consist of a layered set of files. Typically only the top
layer changes when new app versions are released, so network traffic and
storage requirements are smaller than when using monolithic VM images.

VM apps normally include the VM or Docker image as part of the app
version, but it's possible to have it be an input file of the job, in
which case a single (BOINC) app can be used for arbitrarily many science
applications. This is significant because creating an app includes
signing its files, which is a manual operation. We call such an app a
universal app. If a project's science applications are already packaged
as VM or Docker images, they can be run under BOINC with no porting or
other work.

There is a security consideration in the use of universal apps, since
job input files are not code-signed (see Section 3.8). If a project's
server is hacked, the hackers could substitute a different VM image.
However, the danger is small since the VM is sandboxed and can't access
files on the host other than its own input and output files.

VM apps were proposed in 2007 by González et al.~{[}17{]}, and VM
wrappers were developed at CERN {[}18{]} and SZTAKI {[}19{]}. Ferreira
et al.~describe a wrapper able to interface to multiple VM hypervisors
{[}20{]}.

\section{3.7 Job Submission and File
Management}\label{job-submission-and-file-management}

Job submission consists of three steps: 1) staging the job's input files
to a publicly-accessible web server, 2) creating the job, and 3)
handling the results of a

completed job. BOINC supplies C++ APIs and command-line tools for
staging files and submitting jobs on the server. It also provides HTTP
RPC APIs for remote file management and job submission, and it supplies
bindings of these interfaces in Python, PHP, and C++. These interfaces
are defined in terms of ``batches'' of jobs, and are designed for
efficiency: for example, submitting a batch of a thousand jobs takes
less than a second.

The handling of completed jobs is done by a per-app assimilator daemon.
This is linked with a project-supplied C++ or Python function that
handles a completed job. This function might, for example, move the
output files to a different directory, or parse the output files and
write the results to a database.

These mechanisms make it straightforward to use BOINC as a back end for
other job processing systems. For example, we implemented a system where
HTCondor jobs can be automatically routed to a BOINC project. The
mechanisms can also be used for web-based remote job submission systems.
For example, a science gateway might automatically route jobs, submitted
via the web, to either local compute nodes or a BOINC server (see
Section 10.2).

In some BOINC projects, multiple job submitters contend for the
project's resources. BOINC provides an allocation system that handles
this contention in a way that's fair to submitters both with sporadic
and continuous workloads. For each submitter, the system maintains a
balance that grows linearly at a particular rate, up to a fixed maximum.
The rate may differ between submitters; this determines the average
computing throughput available to each submitter. When a submitter uses
resources, their balance is decreased accordingly. At any given point,
the jobs of the submitter with the greatest balance are given priority.
We call this the linear-bounded model; BOINC uses it in other contexts
as well (Sections 6.1 and 10.1). Given a mix of continuous and sporadic
workloads, this policy prioritizes small batches, thereby minimizing
average batch turnaround.

\section{3.8 Storage}\label{storage}

The BOINC storage model is based on named files. Files are immutable: a
file with a given name cannot change. Projects must enforce this. Input
and output files also have logical names, by which applications refer to
them; these need not be unique.

The entities involved in computation -- app versions and jobs -- consist
of collections of files. The files in an app version must be code-signed
by tagging them with a PKE-encrypted hash of the file contents. The
recommended way of doing this is to keep the private key on a computer
that is never connected to a network and is physically secure. Files are
manually copied to this computer -- say, on a USB drive -- and signed.
This procedure provides a critical level of security: hackers are unable
to use BOINC to run malware on volunteer hosts, even if they break into
project servers.

Job input and output files normally are deleted from the volunteer host
after the job is finished. They can be marked as sticky, in which case
they are not deleted. Scheduler requests include a list of sticky files
on the client, and replies can direct the client to delete particular
sticky files.

Files making up app versions are automatically sticky in the sense that
they are deleted only when the app version has been superseded by
another app version with the same app, platform, and plan class, and
with a greater version number.

The amount of storage that BOINC may use on a given host is limited by
its free disk space, and also by volunteer preferences (which can
specify, e.g., minimum free space or maximum fraction of space used by
BOINC). Scheduler requests include the amount of usable space, and the
scheduler doesn't send jobs whose projected disk usage would exceed
this. If the amount of usable space is negative (i.e.~BOINC has exceeded
its limit) the scheduler returns a list of sticky files to delete.

If the host is attached to multiple projects and the disk limit is
reached, the projects are contending for space. The client computes a
storage share for each project, based on its resource share and the disk
usage of the projects; this is what is conveyed to the scheduler.

Files need not be associated with jobs. A scheduler reply can include
lists of arbitrary files to be downloaded, uploaded, or deleted. Thus
BOINC can be used for volunteer storage as well as computing. This might
be used for a variety of purposes: for example, to cache recent data
from a radio telescope so that, if an event is detected a few hours
after its occurrence, the raw data would still be available for
reanalysis. It can also be used for distributed data archival; see
Section 10.3.

\section{3.9 The Anonymous Platform
Mechanism}\label{the-anonymous-platform-mechanism}

By default, the BOINC client gets application executables from the
project server. This model doesn't address the following scenarios:

\begin{itemize}
\tightlist
\item
  The volunteer's computer type is not supported by the project.\\
\item
  The volunteer wants, for security reasons, to run only executables
  they have compiled from source.\\
\item
  The volunteer wants to optimize applications for particular CPU
  models, or to make versions for GPUs or other coprocessors.
\end{itemize}

To handle these cases, BOINC offers a mechanism called anonymous
platform. This lets volunteers build applications themselves, or obtain
them from a third party, rather than getting them from the project
server. This mechanism can be used only for projects that make their
application source code available.

After building their own version of project's applications, the
volunteer creates a configuration file describing these versions: their
files, their CPU and GPU usage, and their expected FLOPS. If the BOINC
client finds such a configuration file, it conveys it to the scheduler,
which uses that set of app versions in the job dispatch process.
Projects that support anonymous platform typically supply input and
output files for a set of ``reference jobs'' so that volunteers can
check the correctness of their versions.

In some projects (for example, SETI@home) the anonymous platform
mechanism has enabled a community of volunteers who have made versions
for various GPUs and vendor-specific CPU features, and who have made
improvements in the base applications. Many of these versions have later
been adopted by the project itself.

\section{3.10 Specialized Job Processing
Features}\label{specialized-job-processing-features}

App versions are assigned version numbers that increase over time. In
general, BOINC always uses the latest version for a given (platform,
plan class) combination. In some cases a project may release a new app
version whose output is different from, and won't validate against, the
previous version. To handle this, BOINC allows jobs to be ``pinned'' to
a particular version number, in which case they will be processed only
by app versions with that version number.

Some BOINC projects have jobs that run for weeks or months on typical
computers {[}21{]}. Because of host churn or user impatience, some of
these jobs don't run to completion. The intermediate outputs of these
jobs may have scientific value. BOINC has two features to support such
jobs. First, applications can tell the client that a particular output
file is complete and can be uploaded immediately, rather than waiting
for the job to finish. Second, the application can generate trickle-up
messages that are conveyed immediately to the server and handled by
project-specific logic. This can be used, for example, to give
volunteers credit for partial job completion.

Some applications have large input files, each of which is used by many
jobs. To minimize network traffic and server load for these
applications, BOINC has a feature called locality scheduling. The large
input files are marked as sticky (see Section 3.8) so that they remain
resident on the host. The scheduler RPC request message includes a list
of sticky files on the host. The scheduler preferentially sends jobs
that use these files (see Section 6.4).

BOINC allows jobs to be targeted at a particular host or volunteer. This
can be used, for example to ensure that test jobs are executed on the
project's own computers.

The difference in performance between volunteer devices -- say, a phone
versus a desktop with a high-end GPU -- can be several orders of
magnitude. A job that completes in 5 min on one could take a week on the
other. To reduce server load and maintain volunteer interest, it's
desirable that job runtime be on the order of an hour. To support
uniformity of job runtime across platforms, BOINC provides a mechanism
that lets projects create jobs in a range of sizes (say, processing
smaller or larger amounts of data). The BOINC server maintains
statistics of the performance of each (host, app version) pair and
dispatches jobs to appropriate hosts based on their size -- larger jobs
are sent to faster hosts.

Applications that use negligible CPU time (e.g., sensor monitoring,
network performance testing, and web crawling) can be labeled as
non-CPU-intensive. BOINC treats these applications specially: it runs a
single job per client, it always runs the application, and it runs it at
normal process priority.

Debugging BOINC applications -- e.g.~figuring out why they crash -- is
challenging because the errors occur on remote, inaccessible machines.
BOINC provides features to assist in this. The client reports exit codes
and

signal numbers, as well as standard error output, to the server. The
server provides web interfaces allowing developers to break down errors
by platform, app version, and error type, and to see the details of
devices where errors occurred; this facilitates finding problems
specific to a particular OS version or video device driver.

\section{4 The Life of a Job}\label{the-life-of-a-job}

In most batch systems the life of a job is simple: it gets dispatched to
a node and executed. In BOINC, however, a dispatched job has various
possible outcomes:

\begin{enumerate}
\def\labelenumi{\arabic{enumi}.}
\tightlist
\item
  The job completes successfully and correctly.\\
\item
  The program crashes.\\
\item
  The job completes but the returned results are incorrect, either
  because of hardware malfunction or because a malicious volunteer
  intentionally substitutes wrong results.\\
\item
  The job is never returned, because (for example) the computer stops
  running BOINC.
\end{enumerate}

To ensure the eventual completion of a job, and to validate its results,
it may be necessary to create additional instances of the job.

Each job J is assigned a parameter \(delay\_bound(J)\) by its submitter.
If an instance of J is dispatched to a host at time T, the deadline for
its completion is \(T + delay\_bound(J)\) . If the completed results
have not been returned to the server by the deadline, the instance is
assumed to be lost (or excessively late) and a new instance of J is
created. The scheduler tries to dispatch jobs only to hosts that are
likely, based on runtime estimates and existing queue sizes, to complete
them by their deadlines.

Projects may have time-sensitive workloads: e.g., batches of jobs that
should be finished quickly because the submitter needs the results to
proceed. Such jobs may be assigned low delay bounds; however, this may
limit the set of hosts to which the job can be dispatched, which may
have the opposite of the desired effect.

For workloads that are purely throughput-oriented, \(delay\_bound(J)\)
can be large, allowing even slow hosts to successfully finish jobs.
However, it can't be arbitrarily large: while the job is in progress,
its input and output files take up space on the server, and
delay\_bound(J) limits this period. For such workloads, a delay bound of
a few weeks may be appropriate.

In addition to \(delay\_bound(J)\) , a job has the following parameters;
these are assigned by the submitter, typically at the level of app
rather than job:

\begin{itemize}
\tightlist
\item
  min\_quorum(J): validation is done when this many successful instances
  have completed. If a strict majority of these instance have equivalent
  results, one of them is selected as the canonical instance, and that
  instance is considered the correct result of the job.\\
\item
  init\_ninstances(J): the number of instances to create initially. This
  must be at least min\_quorum(J); it may be more in order to achieve a
  quorum faster.\\
\item
  max\_error\_instances(J): if the number of failed instances exceeds
  this, the job as a whole fails and no more instances are created. This
  protects against jobs that crash the application.\\
\item
  max\_success\_instances(J): if the number of successful instances
  exceeds this and no canonical instance has been found, the job fails.
  This protects against jobs that produce nondeterministic results.
\end{itemize}

So the life of a job J is as follows:

\begin{itemize}
\tightlist
\item
  \(J\) is submitted, and init\_ninstances(J) instances are created and
  eventually dispatched.\\
\item
  When an instance's deadline is reached and it hasn't been reported,
  the server creates a new instance.\\
\item
  When a successful instance is reported and there is already a
  canonical instance, the new instance is validated against it to
  determine whether to grant credit (see Section 7). Otherwise, if there
  are at least min\_quorum(J) successful instances, the validation
  process is done on these instances. If a canonical instance is found,
  the app's assimilator is invoked to handle it. If not, and the number
  of successful instances exceeds max\_success\_instances(J), J is
  marked as failing.\\
\item
  When a failed instance is reported, and the number of failed instances
  exceeds max\_error\_instances(J), J is marked as failing. Otherwise a
  new instance is created.
\end{itemize}

When a canonical instance is found, J's input files and the output files
of other instances can be deleted, and any unsent instances are
cancelled.

It's possible that a successful instance of \(J\) is reported after a
canonical instance is found. It's important to grant credit to the
volunteer for the instance, but only if it's correct. Hence the output
files of the canonical instance are retained until all instances are
resolved.

After all instances are resolved, and the job has succeeded or failed,
the job and job instance records can be purged from the database. This
keeps the database from growing without limit over time; it serves as a
cache of jobs in progress, not an archive of all jobs. However, the
purging is typically delayed for a few days so that volunteers can view
their completed jobs on the web.

\section{5 Software Architecture}\label{software-architecture}

\section{5.1 Server Architecture}\label{server-architecture}

A BOINC project operates a server consisting of one or more computer
systems running the BOINC server software. These systems may run on
dedicated hardware, VMs, or cloud nodes. The server has many interacting
processes, as shown in Fig. 3.

A BOINC server centers on a relational database (usually MySQL or
MariaDB). The database includes tables for most of the abstractions
mentioned earlier: volunteer accounts, hosts, apps, app versions, jobs,
job instances, and so on. To reduce DB load, details such as the list of
a job's input files are stored in XML ``blobs'' rather than in separate
tables.

Client RPCs are handled by a scheduler, which is implemented as a CGI or
FCGI program run from a web server. Typically many RPCs are in progress
at once, so many instances of the scheduler are active. The main
function of the scheduler is to dispatch job instances to clients, which
requires scanning many jobs and job instances. Doing this by directly
querying the database would limit performance. Instead, BOINC uses an
architecture in which a shared-memory segment contains a cache of job
instances available for dispatch -- typically on the order of a thousand
jobs. This cache is replenished by a feeder process, which fetches
unsent job instances from the database and puts them in vacant slots in
the cache. The feeder supports the homogeneous redundancy, homogeneous
app version, and job size features by ensuring that the shared-memory
cache contains job of all the different categories.

An instance of the scheduler, rather than querying the database for jobs
to send, scans the job cache. It acquires a mutex while doing this; the
scan is fast so there is little contention for the mutex. When the
scheduler sends a job, it clears that entry in the cache, and the entry
is eventually refilled by the feeder. The scheduler updates the database
to mark the job instance as sent. The efficiency of this mechanism
allows a BOINC server -- even on a single machine -- to dispatch
hundreds of jobs per second {[}22{]}.

The server's other functions are divided among a set of daemon
processes, which communicate and synchronize through the database.

Fig. 3 The components of a BOINC server\\
\textit{[Image omitted: images/6a789dfd2fbc5f758c3ea3e79a0382e4776c89960ba8a6e73a5321b8bd5b0972.jpg]}

flowchart

\begin{Shaded}
\begin{Highlighting}[]
\NormalTok{graph TD}
\NormalTok{    A["BOINC client"] {-}{-}\textgreater{} B["scheduler"]}
\NormalTok{    B {-}{-}\textgreater{} C["shared{-}memory job cache"]}
\NormalTok{    C {-}{-}\textgreater{} D["feeder"]}
\NormalTok{    C {-}{-}\textgreater{} E["database"]}
\NormalTok{    E {-}{-}\textgreater{} F["transitioner"]}
\NormalTok{    E {-}{-}\textgreater{} G["DB purger"]}
\NormalTok{    E {-}{-}\textgreater{} H["file deleter"]}
\NormalTok{    F {-}{-}\textgreater{} I["assimilators"]}
\NormalTok{    G {-}{-}\textgreater{} J["validators"]}
\NormalTok{    H {-}{-}\textgreater{} K["job creation"]}
\NormalTok{    C {-}{-}\textgreater{} L["job 0"]}
\NormalTok{    C {-}{-}\textgreater{} M["job 1"]}
\end{Highlighting}
\end{Shaded}

\begin{itemize}
\tightlist
\item
  Validator: There is one for each app that uses replication. It
  examines the instances of a job, compares their output files, decides
  whether there is a quorum of equivalent results, and if so designates
  one instance as ``canonical''. For apps that use homogeneous
  redundancy to achieve bitwise agreement between instances, BOINC
  supplies a validator that compares the output files byte for byte.
  Other apps use custom validators with a project-supplied function that
  does the (fuzzy) comparison for that app.\\
\item
  Assimilator: This handles completed and validated jobs. There is one
  assimilator per app; it includes a project-supplied handler function.
  This function might move the output files to a particular destination,
  or parse them and write the results to a database.\\
\item
  Transitioner: Viewing the progress of a job as a finite-state machine,
  this handles the transitions. The events that trigger transitions come
  from potentially concurrent processes like schedulers and validators.
  Instead of handling the transitions, these programs set a flag in the
  job's database record. The transitioner enumerates these records and
  processes them. This eliminates the need for concurrency control of DB
  access.\\
\item
  File deleter: This deletes input and output files of completed and
  assimilated jobs.\\
\item
  Database purger: This deletes the database records of completed jobs
  and job instances.
\end{itemize}

This multi-daemon architecture provides a form of fault-tolerance. If a
daemon fails (say, an assimilator fails because an external database is
down) the other components continue; work for the failed component
accumulates in the database and is eventually processed.

A BOINC server is highly scalable. First, the daemons can be run on
separate hosts. Second, all of the server functions except the database
server can be divided among multiple processes, running on the same or
on different hosts, by partitioning the space of database job IDs: if
there are N instances of a daemon, instance i handles rows for which
\((ID \mod N) = i\) . In combination with the various options for
scaling database servers, this allows BOINC servers to handle workload
on order of millions of jobs per day.

The BOINC server software also includes scripts and web interfaces for
creating projects, stopping/starting projects, and updating apps and app
versions.

\section{5.2 Client Architecture}\label{client-architecture}

The BOINC client software consists of three programs, which communicate
via RPCs over a TCP connection:

\begin{itemize}
\tightlist
\item
  The core client manages job execution and file transfers. The BOINC
  installer arranges for the client to run when the computer starts up.
  It exports a set of RPCs to the GUI and the screensaver.\\
\item
  A GUI (the BOINC Manager) lets volunteers control and monitor
  computation. They can, for example, suspend and resume CPU and GPU
  computation, and control individual jobs. They can view the progress
  and estimated remaining time of jobs, and can view an event log
  showing a configurable range of internal and debugging information.\\
\item
  An optional screensaver shows application graphics (Section 3.5) for
  running jobs when the computer is idle.
\end{itemize}

Third party developers have implemented other GUIs; for example,
BoincTasks is a GUI for Windows that can manage multiple clients
{[}23{]}. For hosts without a display, BOINC supplies a command-line
program providing the same functions as the Manager.

The core client is multi-threaded. The main thread uses the POSIX
select() call to multiplex concurrent activities (network communication
and time-consuming file operations like decompression and checksumming).
Each GUI RPC connection is serviced by a separate thread. All the
threads synchronize using a single mutex; the main thread releases this
mutex while it's waiting for I/O or doing a file operation. This ensures
rapid response to GUI RPCs.

\section{6 Scheduling}\label{scheduling}

Job scheduling is at the core of BOINC. The scheduling system consists
of three related parts (see Fig. 4):

\begin{itemize}
\tightlist
\item
  The BOINC client maintains a queue of jobs; these jobs use various
  combinations of CPUs and GPUs. Resource scheduling is the decision of
  what jobs to run at a given point.\\
\item
  Work fetch involves deciding when to request more jobs from a project,
  how much work to request, and, if the client is attached to multiple
  projects, which project to ask.
\end{itemize}

\textit{[Image omitted: images/fd617659744b0c08b155b6567d8fc23ca45e2f3fe47ddac39cef8463dc84c6de.jpg]}

flowchart

\begin{Shaded}
\begin{Highlighting}[]
\NormalTok{graph TD}
\NormalTok{    A["Server"] {-}{-}\textgreater{}|Job selection| B["Client"]}
\NormalTok{    B {-}{-}\textgreater{}|Work fetch| A}
\NormalTok{    B {-}{-}\textgreater{} C["Resource scheduling"]}
\NormalTok{    C {-}{-}\textgreater{} D["CPUs"]}
\NormalTok{    C {-}{-}\textgreater{} E["GPUs"]}
\end{Highlighting}
\end{Shaded}

Fig. 4 Job scheduling in BOINC involves three interacting policies:
resource scheduling, work fetch, and job selection

- Job selection is the server's choice of what jobs to send in response
to a client work request.

The scheduling system has several goals:

• Eventually complete all jobs.\\
- Maximize throughput by a) avoiding idle processing resources; b)
completing jobs by their deadline, thereby avoiding creating additional
instances of the job; c) using the best-performing app version for a
given job and host.\\
Enforce the volunteer's per-project resource shares and computing
preferences.

The performance of early versions of the system have been studied using
simulation {[}22, 24, 25, 26{]}. However, the current system has for the
most part been developed empirically, in response to project demands.

BOINC's scheduling policies revolve around estimating the runtime of a
job, on a particular host, using a particular app version. There are
actually two notions of runtime:

\begin{itemize}
\tightlist
\item
  The raw runtime is how long the job takes running full-time on the
  host.\\
\item
  The scaled runtime is how long it takes given a) CPU throttling
  (Section 2.4) and b) the possibly sporadic availability of the
  processing resources used by the app version (resources are
  unavailable
\end{itemize}

when the computer is turned off or their use is disallowed by user
preferences). The client maintains averages of the availability of CPUs
and GPUs, and reports these to the server.

Runtime estimates are intrinsically uncertain. BOINC tries to make good
estimates, and to deal gracefully with situations where the estimates
are inaccurate.

\section{6.1 Client: Resource
Scheduling}\label{client-resource-scheduling}

On a given host, the client manages a set of processing resources: a CPU
and possibly one or more coprocessors such as GPUs. There may be
multiple instances of each processing resource. Volunteer preferences
may limit the number of available CPU instances, denoted
n\_usable\_cpus. The client estimates the peak FLOPS of each instance.
For CPUs this is the Whetstone benchmark result; for GPUs it's the
vendor-supplied estimate based on number of cores and clock rate.

At a given point the client has a queue of jobs to run. A given job may
be a) not started yet; b) currently executing; c) in memory but
suspended, or d) preempted and not in memory.

Each job \(J\) has a fixed resource usage: the number (possibly
fractional) of instances of each processing resource that it uses while
executing. For CPUs, fractional resource usage refers to time: 0.5 means
that \(J\) has 1 thread that runs half the time. For GPUs, fractional
usage typically refers to core usage: 0.5 means that \(J\) uses half the
GPU's cores. In either case, throughput is maximized by running 2 such
jobs at once.

Each job \(J\) has an estimated RAM working set size \(est\_wss(J)\) .
If \(J\) is running or preempted this is the latest working set size
reported by the operating system. If \(J\) is unstarted and there's a
running job \(J_2\) using the same app version, \(est\_wss(J)\) is the
working set size of \(J_2\) ; otherwise it's the project-supplied RAM
usage estimate.

We call a set of jobs feasible if a) for each coprocessor, the sum of
usages by the jobs is at most the number of instances; b) the sum of CPU
usages of CPU jobs is at most n\_usable\_cpus, and the sum of CPU usages
of all jobs is at most n\_usable\_cpus + 1; c) the sum of est\_wss(J) is
at most the amount of available RAM. A set of jobs is maximal if it's
feasible and adding any other queued job would make it infeasible. The
client always runs a maximal set of jobs.

The basic resource scheduling policy is weighted round robin (WRR).
Projects have dynamic scheduling

priorities based on their resource share and recent usage, using the
linear bounded model described in Section 3.7. A maximal set of jobs
from the highest-priority projects are run FIFO. At the end of each time
slice (default 1 h) priorities are recomputed and a possibly different
maximal set of jobs is run. Time-slicing is used so that volunteers see
a mixture of projects.

The WRR policy can result in jobs missing their deadlines. To avoid
this, we need a more sophisticated policy that predicts and avoids
deadline misses. This requires estimating the remaining runtime of jobs.
If a job hasn't started yet this is the job's FLOP estimate (Section
3.2) divided by the server-supplied FLOPs/s estimate (Section 6.3); we
call this the static estimate. If the job has started and has reported a
fraction done (see Section 3.4) this gives a dynamic estimate, namely
the current runtime divided by the fraction done. For some applications
(e.g.~those that do a fixed number of iterations of a deterministic
loop) the dynamic estimate is accurate; projects can flag such
applications, instructing the client to always use the dynamic estimate.
For other applications the fraction done is approximate; in this case
the client uses a weighted average of the static and dynamic estimates,
weighted by the fraction done.

The client periodically simulates the execution of all jobs under the
WRR policy, based on the estimated remaining scaled runtimes of the
jobs. This simulation (described in more detail below) predicts which
jobs will miss their deadlines under WRR.

We can now describe BOINC's resource scheduling policy. The following
steps are done periodically, when a job exits, and when a new job
becomes ready to execute:

\begin{itemize}
\tightlist
\item
  Do the WRR simulation.\\
\item
  Make a list of jobs ordered by the following criteria (in descending
  priority): a) prefer jobs that miss their deadline under WRR
  simulation, and run these
\end{itemize}

earliest deadline first; b) prefer GPU jobs to CPU jobs; c) prefer jobs
in the middle of a time slice, or that haven't checkpointed; d) prefer
jobs that use more CPUs; e) prefer jobs from projects with higher
scheduling priority.

- Scan through this list, adding jobs until a maximal set is found. Run
these jobs, and preempt running jobs not in the set.

In essence, BOINC uses WRR scheduling unless there are projected
deadline misses, in which case it uses earliest deadline first (EDF).
EDF is optimal on uniprocessors but not multiprocessors {[}27{]}; a
policy such as least slack time first might work better in some cases.

\section{6.2 Client: Work Fetch Policy}\label{client-work-fetch-policy}

Volunteer preferences include lower and upper limits, denoted \(B_{LO}\)
and \(B_{HI}\) , on the amount of buffered work, as measured by
remaining estimated scaled runtime. When the amount of work for a given
processing resource falls below \(B_{LO}\) , the client issues a request
to a project, asking for enough work to bring the buffer up to at least
\(B_{HI}\) .

This buffering scheme has two purposes. First, the client should buffer
enough jobs to keep all processing resources busy during periods when no
new jobs can be fetched, either because the host is disconnected from
the Internet, project servers are down, or projects have no jobs.
\(B_{LO}\) should reflect the expected duration of these periods.
Second, the frequency of scheduler RPCs should be minimized in order to
limit load on project servers; an RPC should get many jobs if possible.
The difference \(B_{HI} - B_{LO}\) expresses this granularity.

The WRR simulation described above also serves to estimate buffer
durations; see Fig. 5.

Fig. 5 Example of WRR simulation for a given processing resource. Dark
gray bars represent instance usage. The area of light gray bars is the
resource's ``shortfall''\\
resource instances (e.g.~CPUs)\\
\textit{[Image omitted: images/f82532628a84f4e4b3f7c050f4a5e7d03d6e72c1e89e02251f6ba9e4b2c23afa.jpg]}

text\_image

1 2 3 4 now now + B\_LO now + B\_HI time

In the simulation, each instance A of the processing resource R is used
for some period \(T(A)\) starting at the current time. If
\(T(A) < B_{LO}\) for any instance, then the buffer for that resource
needs to be replenished. \(Shortfall(R)\) denotes the minimum additional
job duration needed to bring \(T(A)\) up to at least \(B_{HI}\) for all
instances A, namely

\[
\operatorname{shortfall} (R) = \sum_ {\text { instances   A   of   } R} \max (0, B _ {H I} - T (A))
\]

A work request includes several components for each processing resource
R:

\begin{itemize}
\tightlist
\item
  req\_runtime(R): the buffer shortfall. The server should send a set of
  jobs whose scaled runtime estimates times the resource usage exceeds
  this.\\
\item
  req\_idle(R): The number of idle instances. The server should send
  jobs whose total usage of R is at least this.\\
\item
  queue\_dur(R): the estimated remaining scaled runtime of queued jobs
  that use R. This is used to estimate the completion time of new jobs
  (Section 6.4).
\end{itemize}

There are situations in which a work request can't be issued to a
project \(P\) : for example, \(P\) is suspended by the user, or the
client is in exponential backoff from a previous failed RPC to \(P\) .
There are also situations where work for \(R\) can't be requested from
\(P\) : for example, if it has no app versions that use \(R\) , or user
preferences prohibit using \(R\) . If none of these conditions hold, we
say that \(R\) is fetchable for \(P\) .

We can now describe the work fetch policy. First, perform the WRR
simulation. If some resource needs to be replenished, scan projects in
order of decreasing scheduling priority (Section 6.1). Find the first
one P for which some processing resource R needs to be replenished, and
R is fetchable for P. Issue a scheduler RPC to P, with work requests
parameters as described above for all processing resources that are
fetchable for P.

Sometimes the client makes a scheduler RPC to a project P for reasons
other than work fetch - for example, to report a completed job instance,
because of a user request, or because an RPC was requested by the
project. When this occurs, the client may ``piggyback'' a work request
onto the RPC. For each processing resource type R, it sees if P is the
project with highest scheduling priority for which R is fetchable. If
so, it requests work for R as described above.

The client doesn't generally report completed jobs immediately; rather,
it defers until several jobs can be reported in one RPC, thus reducing
server load. However, it always reports a completed job when its
deadline is approaching or if the project has requested that results be
reported immediately.

\section{6.3 Server: Job Runtime
Estimation}\label{server-job-runtime-estimation}

We now turn to the server's job selection policy. Job runtime estimation
plays two key roles in this. First, the scheduler uses runtime estimates
to decide how many jobs satisfy the client's request. Second, it uses
the estimated turnaround time (calculated as runtime plus queue
duration) to decide whether the job can be completed within its delay
bound.

The job submitter may have a priori information about job duration; for
example, processing a data file that's twice as large as another might
be known to take twice as long. This information is expressed in a
parameter est\_flop\_count(J) supplied by the job submitter (see Section
3.2). This is an estimate of the number of FLOPs the job will use, or
more precisely TR, where T is the expected runtime of the job on a
processing resource whose peak FLOPs/s is R. If job durations are
unpredictable, est\_flop\_count() can be a constant.

The BOINC server maintains, for each pair of host H and app version V,
the sample mean \(\overline{R}(H, V)\) and variance of
(runtime(J)/est\_flop\_count(J)) over all jobs J processed by H and V.
It also maintains, for each app version V, the sample mean
\(\overline{R}(V)\) of (runtime(J)/est\_flop\_count(J)) for jobs
processed using V across all hosts.

For a given host H and app version V, the server computes a projected
FLOPS proj\_flops(H, V), which can be thought of as the estimated FLOPS,
adjusted for systematic error in est\_flop\_count(). proj\_flops(H, V)
is computed as follows.

\begin{itemize}
\tightlist
\item
  If the number of samples of \(\overline{R}(H, V)\) exceeds a threshold
  (currently 10) then proj\_flops(H, V) is \(\overline{R}(H, V)\) .\\
\item
  Otherwise, if the number of samples of \(\overline{R}(V)\) exceeds a
  threshold, \(proj\_flops(H, V)\) is \(\overline{R}(V)\) .\\
\item
  Otherwise, \(proj\_flops(H, V)\) is the peak FLOPS of \(V\) running on
  \(H\) , i.e.~the sum over processing resources \(R\) of \(V\) 's usage
  of \(R\) times the peak FLOPS of \(R\) .
\end{itemize}

The estimated raw runtime of job J on host H using app version V,
denoted est\_runtime(J, H, V), is then
est\_flop\_count(J)/proj\_flops(H, V). proj\_flops(H, V) is also used to
select which app version to use to process J on H, as explained below.

\section{6.4 Server: Job Dispatch
Policy}\label{server-job-dispatch-policy}

A scheduler RPC request message from a host H includes a description of
H, including its OS, processing resources (CPUs and GPUs), memory and
disk. It also includes, for each processing resource R, the fraction of
time R is available for use, and the request parameters described in
Section 6.2.

The RPC reply message includes a list of jobs, and for each job a) the
app version to use, and b) an estimate of the FLOPS that will be
performed by the program. The general goal of the scheduler is to
dispatch jobs that satisfy the request, and to use the app versions that
will run fastest.

The scheduler loops over processing resources, handling GPUs first so
that they'll be kept busy if the disk space limit is reached (as it
dispatches jobs, the scheduler keeps track of how much disk space is
still available). For each processing resource \(R\) , the scheduler
scans the job cache, starting at a random point to reduce lock conflict,
and creates a candidate list of possible jobs to send, as follows.

For each scanned job \(J\) , we find the app versions that a) use a
platform supported by \(H\) ; b) use \(R\) ; c) pass homogeneous
redundancy and homogeneous app version constraints, and d) pass the plan
class test. If no such app version exists, we continue to the next job.
Otherwise we select the app version \(V\) for which
\(proj\_flops(H, V)\) is greatest.

We then compute a score for J, representing the ``value'' of sending it
to H. The score is the weighted sum of several factors:

\begin{itemize}
\item
  If \(J\) has keywords (see Section 2.4), and the volunteer has keyword
  preferences, whether \(J\) 's keywords match the volunteer's ``yes''
  keywords; skip \(J\) if it has a ``no'' keyword.\\
\item
  The allocation balance of the job submitter (see Section 3.7).\\
\item
  Whether J was skipped in handling a previous scheduler request. The
  idea is that if J is hard to send, e.g.~because of large RAM
  requirements, we should send it while we can.
\item
  If the app uses locality scheduling, whether \(H\) already has the
  input files required by \(J\) (see Section 3.10).\\
\item
  If the app has multi-size jobs, whether proj\_flops(H, V) lies in the
  same quantile as the size of J (see Section 3.10).
\end{itemize}

We then scan the candidate list in order of descending score. For each
job J, we acquire a semaphore that protects access to the job cache, and
check whether another scheduler process has already dispatched J, and
skip it if so.

We then do a ``fast check'' (with no database access) to see whether we
can actually send J to H, given the jobs we have already decided to
send. We skip J if a) the remaining disk space is insufficient; b) J
probably won't complete by its deadline: i.e.~if queue\_dur(R) +
est\_runtime(J, H, V) \textgreater{} delay\_bound(J); or c) we are
already sending another instance of the same job as J.

We then flag J as ``taken'' in the job cache, release the mutex, and do
a ``slow check'' involving database access: Have we already sent an
instance of J to this volunteer? Has J errored out since we first
considered it? Does it still pass the homogeneous redundancy check? If
so, skip J and release it in the cache.

If we get this far, we add the pair \((J, V)\) to the list of jobs to
send to client, and free the slot in the job cache. Letting E be the
estimated scaled runtime of the job, we add E to queue\_dur(R), subtract
E from req\_runtime(R), and subtract J's usage of R from req\_idle(R).
If both of the latter are nonpositive, we move on to the next processing
resource. Otherwise we move to the next job in the candidate list.

\section{7 Credit}\label{credit}

BOINC grants credit -- an estimate of FLOPs performed -- for completed
jobs. The total and exponentially-weighted recent average credit is
maintained for each computer, volunteer, and team. Credit serves several
purposes: it provides a measure of progress for individual volunteers;
it is the basis of competition between volunteers and teams; it provides
a measure of computational throughput for projects and for BOINC as a
whole; and it can serve as proof-of-work for virtual currency systems.

To support these uses, the credit system ideally should:

\begin{itemize}
\tightlist
\item
  be device neutral: similar jobs (in particular the instances of a
  replicated job) should get about the same credit regardless of the
  host and computing resource on which they are processed;\\
\item
  be project neutral: a given device should get about the same credit
  per time regardless of which project it computes for;\\
\item
  resist credit cheating: i.e.~efforts to get credit for jobs without
  actually processing them.
\end{itemize}

The basic formula is that one unit of credit represents 1 day of a CPU
with a Whetstone benchmark of 1 GFLOPS. Credit is assigned in different
ways depending on the properties of the application. If all jobs do the
same computation, the project can time the job on a typical computer
with known Whetstone benchmark, compute the credit accordingly, and tell
BOINC to grant that amount per job. Similarly, if the app consists of a
loop that executes an unpredictable number of times but that always does
the same computation, the credit per iteration can be computed ahead of
time, and credit can be granted (in the validator) based on the number
of loop iterations performed.

For applications not having these properties, a different approach is
needed. The BOINC client computes the peak FLOPS of processing resources
as described earlier. The actual FLOPS of a particular application will
be lower because of memory access, I/O, and synchronization. The
efficiency of the application on a given computer is actual FLOPS
divided by peak FLOPS. The efficiency of a CPU application can vary by a
factor of 2 between computers; hence computing credit based on runtime
and benchmarks would violate device neutrality. The problem is worse for
GPUs, for which efficiencies are typically much less than for CPUs, and
vary more widely.

To handle these situations, BOINC uses an adaptive credit system. For a
completed job instance J, we define its peak FLOP count \(PFC(J)\) as:

\[
P F C (J) = \sum_ {r \in R} \text { runtime } (J) * \text { usage } (r) * \text { peak\_flops } (r)
\]

where R is the set of processing resources used by the app version.

The server maintains, for each (host, app version) pair and each app
version, the statistics of \(PFC(J)/\) est\_flop\_count(J) where
est\_flop\_count(J) is the a priori job size estimate.

The claimed credit for J is \(PFC(J)\) times two normalization factors:

\begin{itemize}
\tightlist
\item
  Version normalization: the ratio of the version's average PFC to that
  of the version for which average PFC is lowest (i.e.~the most
  efficient version).\\
\item
  Host normalization: the ratio of the (host, app version)'s average PFC
  to the average PFC of the app version.
\end{itemize}

The instances of a replicated job may have different claimed credit.
BOINC computes a weighted average of these, using a formula that reduces
the impact of outliers, and grants this amount of credit to all the
instances.

In the early days of BOINC we found that volunteers who had accumulated
credit on one project were reluctant to add new projects, where they'd
be starting with zero credit. To reduce this effect we introduced the
idea of cross-project credit: the sum of a computer's, volunteer's, or
team's credit over all projects in which they participate. This works as
follows:

\begin{itemize}
\tightlist
\item
  Each host is assigned a unique cross-project ID. The host ID is
  propagated among the projects to which the host is attached, and a
  consensus algorithm assigns a single ID.\\
\item
  Each volunteer is assigned a unique cross-project ID. Accounts with
  the same email address are equated. The ID is based on email address,
  but cannot be used to infer it.\\
\item
  Projects export (as XML files) their total and recent average credit
  for hosts and volunteers, keyed by cross-project ID.\\
\item
  Cross-project statistics web sites (developed by 3rd parties)
  periodically read the statistics data from all projects, collate it
  based on cross-project IDs, and display it in various forms (volunteer
  and host leaderboards, project totals, etc.).
\end{itemize}

\section{8 Supporting Volunteers}\label{supporting-volunteers}

BOINC is intended to be usable for computer owners with little technical
knowledge -- people who may have never installed an application on a
computer. Such people often have questions or require help. We've crowd-

sourced technical support by creating systems where experienced
volunteers can answer questions and solve problems for beginners. Many
issues are handled via message boards on the BOINC web site. However,
these are generally in English and don't support real-time interaction.
To address these problems we created a system, based on Internet
telephony, that connects beginners with questions to helpers who speak
the same language and know about their type of computer. They can then
interact via voice or chat.

BOINC volunteers are international -- most countries are represented.
All of the text in the BOINC interface -- both client GUIs and web site
-- is translatable, and we crowd-source the translation. Currently 26
languages are supported.

\section{9 Related Work}\label{related-work}

Luis Sarmenta coined the term ``volunteer computing'' {[}28{]}. The
earliest VC projects (1996-97) involved prime numbers (GIMPS) and
cryptography (distributed.net). The first large-scale scientific
projects were Folding@Home and SETI@home, which launched in 1999.

Several systems have been built that do volunteer computing in a web
browser: volunteers open a particular page and a Java or Javascript
program runs {[}29{]}. These systems haven't been widely used for
various reasons: volunteers eventually close the window, and C and
FORTRAN science applications aren't supported.

In the context of BOINC, researchers have studied the availability of
volunteer computers {[}5, 30{]} and their hardware and software
attributes {[}11{]}.

It is difficult to study BOINC's scheduling and validation mechanisms in
the field. Initially, researchers used simulation to study client
scheduling policies {[}25, 31{]}, server policies {[}32{]}, and entire
VC systems {[}33{]}. As the mechanisms became more complex, it was
difficult to model them accurately, and researchers began using
emulation -- simulators using the actual BOINC code to model client and
server behavior (we restructured the code to make this possible).
Estrada et al.~developed a system called EmBoinc that combines a
simulator of a large population of volunteer hosts (driven either by
trace data or by a random model) with an emulator of a project server --
that is, the actual server software modified to take input from the
population simulator, and to use virtual time instead of real time
{[}26{]}. Our group developed an emulator of the client {[}24{]}.
Volunteers experience problems can upload their BOINC state files, and
run simulations, through a web-based interface. This lets us debug
host-specific issues without access to the host.

Much work has been done toward integrating BOINC with other HTC systems.
Working with the HTCondor group, we developed an adaptor that migrates
HTCondor jobs to a BOINC project. This is being used by CERN {[}34{]}.
Other groups have built systems that dynamically assign jobs to
different resources (local, Grid, BOINC) based on job size and load
conditions {[}35{]}. Andrzejak et al.~proposed using volunteer resources
as part of clouds {[}36{]}. The integration of BOINC with the European
Grid Infrastructure (EGEE) is described by Visegradi et al.~{[}37{]} and
Kacsuk et al.~{[}38{]}. Bazinet and Cummings discuss strategies for
subdividing long variable-sized jobs into jobs sized more appropriately
to VC {[}39{]}.

Research has been done on several of the problems underlying volunteer
computing, such as validating results {[}40, 41, 42{]}, job runtime
estimation {[}43{]}, accelerating the completion of batches and DAGs of
jobs {[}44, 45{]}, and optimizing replication policies {[}46, 47{]}.
Many of these ideas could and perhaps will be implemented in BOINC. Cano
et al.~surveyed the security issues inherent in VC {[}48{]}.

Files such as executables must be downloaded to all the computers
attached to a project. This can impose a heavy load on the project's
servers and outgoing Internet connection. To reduce this load, Elwaer et
al.~{[}49{]} proposed using a peer-to-peer file distribution system
(BitTorrent) to distribute such files, and implemented this in a
modified BOINC client. This method hasn't been used by any projects, but
it may be useful in the future.

Researchers have explored how existing distributed computing paradigms
and algorithms can be adapted to VC. Examples include evolutionary
algorithms {[}50, 51{]}, virtual drug discovery {[}52{]}, MPI-style
distributed parallel programming {[}53{]}, Map-Reduce computing
{[}54{]}, data analytics {[}55{]}, machine learning {[}56, 57, 58{]},
gene sequence alignment {[}59{]}, climate modeling {[}21{]}, and
distributed sensor networks {[}60{]}. Desell et al.~describe a system
that combines VC and participatory ``citizen science'' {[}61{]}.

Silva et al.~proposed extending BOINC to allow job submission by
volunteers {[}62{]}. Marosi et al.~proposed a cloud-like system based on
BOINC with extensions developed at SZTAKI {[}19{]}.

Binti et al.~studied a job dispatch policy designed to minimize
air-conditioning costs in campus grids {[}63{]}.

\section{10 Conclusion and Future
Work}\label{conclusion-and-future-work}

Volunteer computing can provide high-throughput computing power in
larger quantities and at lower costs than computing centers and
commercial clouds. BOINC provides a flexible open-source platform for
volunteer computing. Section 1 outlined the challenges that arise in
using untrusted, unreliable, and highly heterogeneous computing
resources for HTC. In BOINC, we have developed mechanisms that
effectively address each of these challenges, thus transforming a global
pool of consumer devices into a reliable computing resource for a wide
range of scientific applications.

In addition to summarizing the features, architecture, and
implementation of BOINC, we have presented detailed descriptions of the
algorithms and policies it uses, some of which are novel. These
mechanisms have been developed in response to the demands of real-world,
operational VC projects. They have satisfied these demands, thus showing
their effectiveness. We have not, for the most part, formally studied
the performance of the mechanisms, or compared them with alternatives.
We hope that this paper enables studies of this sort.

We now describe current and prospective future work aimed at extending
the use of BOINC.

\section{10.1 Coordinated Volunteer
Computing}\label{coordinated-volunteer-computing}

The original BOINC participation model was intended to produce a dynamic
and growing ecosystem of projects and volunteers. This has not happened:
the set of projects has been mostly static, and the volunteer population
has gradually declined. The reasons are likely inherent in the model.
Creating a BOINC project is risky: it's a significant investment, with
no guarantee of any volunteers, and hence of any computing power.
Publicizing VC in general is difficult because each project is a
separate brand, presenting a diluted and confusing image to the public.
Volunteers tend to stick with the same projects, so it's difficult for
new projects to get volunteers.

To address these problems, we created a new coordinated model for VC.
This involves an account manager, Science United (SU), in which
volunteers register for scientific areas (using the keyword mechanism
described in Section 2.4) rather than for specific projects {[}64{]}. SU
dynamically attaches hosts to projects based on these science
preferences. Projects are vetted by SU, and SU has a mechanism (based on
the linear-bounded model described in Section 3.7) for allocating
computing power among projects. This means that a prospective new
project can be guaranteed a certain amount of computing power before any
investment is made, thus reducing the risk involved in creating a new
project.

\section{10.2 Integration with Existing HTC
Systems}\label{integration-with-existing-htc-systems}

While the coordinated model reduces barriers to project creation, the
barriers are still too high for most scientists. To address this, we are
exploring the use of VC by existing HTC providers like computing centers
and science gateways. Selected jobs in such systems will be handled by
VC, thus increasing the capacity available to the HTC provider's
clients, which might number thousands of scientists. We are prototyping
this type of facility at Texas Advanced Computing Center (TACC) and at
nanoHUB, a gateway for nanoscience {[}65{]}. Both projects use Docker
and a universal VM app to handle application porting.

\section{10.3 Data Archival on Volunteer
Devices}\label{data-archival-on-volunteer-devices}

Host churn means that a file stored on a host may disappear forever at
any moment. This makes it challenging to use BOINC for high-reliability
data archival. However, it's possible. Encoding (e.g.~Reed-Solomon) can
be used to store a file in a way that allows it to be reconstructed even
if a number of hosts disappear. This can be used to achieve arbitrarily
high levels of reliability with reasonable space overhead. However,
reconstructing the file requires uploading the remaining chunks to a
server. This imposes high network overhead, and it requires that the
original file fit on the server. To solve these problems, we have
developed and implemented a technique called ``multi-level encoding'' in
which the chunks from a top-level encoding are further encoded into
smaller 2nd-level chunks. When a failure occurs, generally only a
top-level chunk needs to be reconstructed, rather than the entire file.

\section{10.4 Using More Types of
Devices}\label{using-more-types-of-devices}

Ideally, BOINC should be available for any consumer device that's
networked and can run scientific jobs. This is not yet the case. BOINC
runs on Android but not Apple iOS, because Apple's rules bar apps from
dynamically downloading executable code. It also is not available for
video game consoles or Internet of Things devices such as appliances.
The barriers are generally not technical, but rather that the
cooperation of the device manufacturers is needed.

In addition, although BOINC on Android detects on-chip GPUs via OpenCL,
no projects currently have apps that use them. It's likely that their
use will raise issues involving heat and power.

\section{10.5 Secure Computing}\label{secure-computing}

Recent processor architectures include features that allow programs to
run in secure ``enclaves'' that cannot be accessed by the computer
owner. Examples include Intel's Software Guard Extensions (SGX), AMD
Secure Encrypted Virtualization, ARM's TrustZone, and the proposed
Trusted Platform Module {[}66{]}. Extending BOINC to use such systems
has two possible advantages: a) it could eliminate the need for result
validation, and b) it would make it possible to run applications whose
data is sensitive -- for example, commercial drug discovery, and medical
research applications involving patient data.

\section{10.6 Recruiting and Retaining
Volunteers}\label{recruiting-and-retaining-volunteers}

Early volunteer computing projects such as SETI@home received mass media
coverage in 1999-2001, and attracted on the order of 1 M volunteers.
Since then there has been little media coverage, and the user base has
shrunk to 200 K or so. Surveys have shown that most volunteers are male,
middle-aged, IT professionals {[}67, 68{]}.

Several technology companies have promoted VC by developing and
publicizing ``branded'' versions of the BOINC client. Examples include
Intel (Progress thru Processors), HTC (Power to Give) and Samsung (Power
Sleep). None of these resulted in a significant long-term increase in
volunteers.

Volunteer computing faces the challenge of marketing itself to the
broader public. Possible approaches include social media, public service
announcements, and partnerships with companies such as video game
frameworks, device manufacturers, and cell phone service providers.

Another idea is to use BOINC-based computing as proof-of-work for
virtual currency and other blockchain systems, and use the tokens issued
by these systems as a reward for participating. One such system exists
{[}69{]}, and others are under development.

\section{10.7 Scheduling}\label{scheduling-1}

There are a number of possible extensions and improvements related to
scheduling; some of these are explored by Chernov et al.~{[}70{]}. Dinis
et al.~proposed facilitating experimentation by making the server
scheduler ``pluggable'' {[}71{]}.

BOINC currently has no direct support for directed acyclic graphs (DAGs)
of dependent jobs. BOINC doesn't currently support low-latency
computing, by which we mean the ability to complete a job or DAG of jobs
as fast as possible, and to give an estimate of completion time in
advance. Doing so would require a reworking of all the scheduling
mechanisms. For example, job runtime estimation would need to use sample
variance as well as sample mean, in order to bound the probability of
deadline miss. Job deadlines would be assigned dynamically during job
dispatch, rather than specified by the submitter. Replication could be
used to reduce expected latency.

When using host availability as part of job runtime estimation, BOINC
uses the average availability. This doesn't take into account the
possibility that unavailability occurs on large chunks, or that it
occurs with a daily or weekly periodicity. Doing so could provide more
accurate, and in some cases lower, estimates.

To accelerate the completion of batches (and perhaps DAGs), BOINC could
add a mechanism to deal with the straggler effect: the tendency for a
few jobs to delay the completion of the batch as a whole. The idea is
send these jobs to fast, reliable, and available computers, and possibly
to replicate the jobs. These ideas have been explored by others
{[}44{]}.

Although individual volunteered computers are unpredictable, the
behavior of a large pool of such computers is fairly predictable. Hence
BOINC could potentially provide quality-of-service guarantees, such as a
level of throughput over a time interval, to projects or job submitters
within a project.

Acknowledgements Many people contributed to the ideas presented here and
to their implementation. Thanks in particular to Bruce Allen, Nicolás
Alvarez, Matt Blumberg, Karl Chen, Carl Christensen, Charlie Fenton,
David Gedye, Francois Grey, Eric Korpela, Janus Kristensen, John McLeod
VII, Kevin Reed, Ben Segal, Michela Taufer, and Rom Walton. The National
Science Foundation has supported BOINC with several grants, most
recently award \#1664190.

Open Access This article is distributed under the terms of the Creative
Commons Attribution 4.0 International License
(http://creativecommons.org/licenses/by/4.0/), which permits
unrestricted use, distribution, and reproduction in any medium, provided
you give appropriate credit to the original author(s) and the source,
provide a link to the Creative Commons license, and indicate if changes
were made.

\section{References}\label{references}

\begin{enumerate}
\def\labelenumi{\arabic{enumi}.}
\item
  ``Publications by BOINC projects,''
  https://boinc.berkeley.edu/wiki/Publications\_by\_BOINC\_projects,
  (2018)\\
\item
  Reference.com, ``How Many Computers are There in the World?,''
  https://www.reference.com/technology/many-computers-world-e2e980daa5e128d0,
  (2014)\\
\item
  iTers News, ``Huawei to launch Kirin 980,''
  http://itersnews.com/?p=107208, (2018)\\
\item
  Toth, D.: ``Increasing Participation in Volunteer Computing,'' in IEEE
  Parallel and Distributed Processing Workshops and PhD Forum (IPDPSW),
  Shanghai, (2011)\\
\item
  Heien, E., Kondo, D., Anderson, D.P.: A correlated resource model of
  internet end hosts. IEEE Trans. Parallel Distribut. Syst. 23(6),
  977--984 (2012)\\
\item
  Mohammadi, M., Bazhirov, T.: ``Comparative benchmarking of cloud
  computing vendors with high performance linpack,'' in Proceedings of
  the 2nd International Conference on High Performance Compilation,
  Computing and Communications, Hong Kong, (2018)\\
\item
  Amazon, ``Amazon EC2 Pricing,''
  https://aws.amazon.com/ec2/pricing/on-demand/, (2019)\\
\item
  Amazon, ``Amazon EC2 Spot Instances Pricing,''
  https://aws.amazon.com/ec2/spot/pricing/, (2019)\\
\item
  Wikipedia, ``Nvidia Tesla,''
  https://en.wikipedia.org/wiki/Nvidia\_Tesla, (2019)\\
\item
  Kondo, D., Bahman, J., Malecot, P., Cappello, F., Anderson, D.:
  ``Cost-Benefit Analysis of Cloud Computing versus Desktop Grids,'' in
  18th International Heterogeneity in Computing Workshop, Rome, (2009)\\
\item
  Anderson, D. P., Reed, K.: ``Celebrating diversity in volunteer
  computing,'' in 42nd Hawaii International Conference on System
  Sciences (HICSS '09), Waikoloa, HI, (2009)\\
\item
  Anderson, D. P.: ``BOINC: A System for Public-Resource Computing and
  Storage,'' in 5th IEEE/ACM International Workshop on Grid Computing,
  Pittsburgh, PA, (2004)\\
\item
  McIntosh, E., Schmidt, F., Dinechin, F. D.: ``Massive Tracking on
  Heterogeneous Platforms,'' in Proc. of 9th International Computational
  Accelerator Physics Conference, (2006)
\item
  Taufer, M., Anderson, D., Cicotti, P., III, C. B.: ``Homogeneous
  Redundancy: a Technique to Ensure Integrity of Molecular Simulation
  Results Using Public Computing,'' in 19th IEEE International Parallel
  and Distributed Processing Symposium (IPDPS'05) Heterogeneous
  Computing Workshop, Denver, (2005)\\
\item
  Marosi, A. C., Balaton, Z., Kacsuk, P.: ``GenWrapper: A generic
  wrapper for running legacy applications on desktop grids,'' in 2009
  IEEE International Symposium on Parallel \& Distributed Processing,
  Rome, (2009)\\
\item
  Merkel, D.: Docker: lightweight Linux containers for consistent
  development and deployment. Linux J. 2014, 239 (2014)\\
\item
  González, D. L., Vega, F. F. D., Trujillo, L., Olague, G., Segal, A.
  B.: ``Customizable Execution Environments with Virtual Desktop Grid
  Computing,'' in Parallel and Distributed Computing and Systems (PDCS
  2007), Cambridge, MA, (2007)\\
\item
  Segal B., Buncic P., Quintas, D., Gonzalez, D., Harutyunyan, A.,
  Rantala, J., Weir, D.: ``Building a volunteer cloud,'' in Conferencia
  Latinoamericana de Computación de Alto Rendimiento, Mérida, Venezuela,
  2009\\
\item
  Marosi, A., Kovacs, J., Kacsuk, P.: Towards a volunteer cloud system.
  Futur. Gener. Comput. Syst. 29, 1442--1451 (2014)\\
\item
  Ferreira, D., Araujo, F., Domingues, P.: ``libboincexec: A Generic
  Virtualization Approach for the BOINC Middleware,'' in 2011 IEEE
  International Symposium on Parallel and Distributed Processing
  Workshops and PhD Forum (IPDPSW), Anchorage, (2011)\\
\item
  Christensen, C., Aina, T., Stainforth, D.: ``The Challenge of
  Volunteer Computing With Lengthy Climate Model Simulations,'' in First
  IEEE International Conference on e-Science and Grid Technologies,
  Melbourne, (2005)\\
\item
  Anderson, D. P., Korpela, E., Walton, R.: ``High-Performance Task
  Distribution for Volunteer Computing,'' in First IEEE International
  Conference on e-Science and Grid Technologies, Melbourne, (2005)\\
\item
  Efmer, F.: ``Boinc Tasks,'' https://efmer.com/boinctasks/, (2018)\\
\item
  Anderson, D. P.: ``Emulating Volunteer Computing Scheduling
  Policies,'' in Fifth Workshop on Desktop Grids and Volunteer Computing
  Systems (PCGrid 2011), Anchorage, (2011)\\
\item
  Kondo, D., Anderson, D. P., McLeod, J.: ``Performance Evaluation of
  Scheduling Policies for Volunteer Computing,'' in 3rd IEEE
  International Conference on e-Science and Grid Computing, Banagalore,
  India, (2007)\\
\item
  Estrada, T., Taufer, M., Anderson, D.: ``Performance prediction and
  analysis of BOINC projects: an empirical study with EmBOINC. J. Grid
  Comput. vol.~7, no. 537, (2007)\\
\item
  Liu, C., Layland, J.: Scheduling algorithms for multiprogramming in a
  hard-real-time environment. J. ACM. 20(1), 46--61 (1973)\\
\item
  Sarmenta, L., Hirano, S., Ward, S.: ``Towards Bayanihan: building an
  extensible framework for volunteer computing using Java''.
  Concurrency: Pract. Exp. vol.~10, no 11--13 (1998)
\item
  Fabisiak, T., Danilecki, A.: ``Browser-based harnessing of voluntary
  computational power''. Found. Comput. Decis. Sci. J. Poznan Univ.
  Technol. vol 42, no 1 (2017)\\
\item
  Javadi, B., Kondo, D., Vincent, J.-M., Anderson, D.P.: Discovering
  statistical models of availability in large distributed systems: an
  empirical study of SETI@home. IEEE Trans. Parallel Distribut. Syst.
  22(11), 1896--1903 (2011)\\
\item
  Toth, D., Finkel, D.: ``Improving the productivity of volunteer
  computing by using the most effective task retrieval policies''. J.
  Grid Comp. vol 7, no 519 (2009)\\
\item
  Taufer, M., Kerstens, A., Estrada, T., Flores, D., Teller, P.:
  ``SimBA: a Discrete Event Simulator for Performance Prediction of
  Volunteer Computing Projects,'' in Proceedings of the International
  Workshop on Principles of Advanced and Distributed Simulation, San
  Diego (2007)\\
\item
  Donassolo, B., Casanova, H., Legrand, A., Velho, P.: ``Fast and
  scalable simulation of volunteer computing systems using SimGrid,'' in
  Proceedings of the 19th ACM International Symposium on High
  Performance Distributed Computing, Chicago (2010)\\
\item
  Barranco, J., Cai, Y., Cameron, D., et al.: LHC@home: a BOINC-based
  volunteer computing infrastructure for physics studies at CERN. Open
  Eng. 7(1), 379--393 (2017)\\
\item
  Myers, D. S., Bazinet, A. L., Cummings, M. P.: ``Expanding the reach
  of Grid computing: combining Globus- and BOINC-based systems,'' in
  Grids for Bioinformatics and Computational Biology, Wiley Book Series
  on Parallel and Distributed Computing, New York, John Wiley \& Sons pp
  71--85 (2008)\\
\item
  Andrzejak, A., Kondo, D., Anderson, D. P.: ``Exploiting Non-Dedicated
  Resources for Cloud Computing,'' in 12th IEEE/IFIP Network Operations
  \& Management Symposium, Osaka, Japan (2010)\\
\item
  Visegradi, A., Kovacs, J., Kacsuk, P.: Efficient extension of gLite
  VOs with BOINC based desktop grids. Futur. Gener. Comput. Syst. 32,
  13--23 (2014)\\
\item
  Kacsuk, P., Kovacs, J., Farkas, A., Marosi, C., Balaton, Z.: Towards a
  powerful European DCI based on desktop grids. J. Grid Comput. 9(2),
  219--239 (2011)\\
\item
  Bazinet, A.L., Cummings, M.P.: Subdividing long-running,
  variable-length analyses into short, fixed-length BOINC Workunits. J.
  Grid Comput. 7(4), 501--518 (2009)\\
\item
  Golle, P., Mironov, I.: ``Uncheatable Distributed Computations,'' in
  Proceedings of the 2001 Conference on Topics in Cryptology: The
  Cryptographer's Track at RSA. (2001)\\
\item
  Sonnek, J., Chandra, A., Weissman, J.: ``Adaptive reputation-based
  scheduling on unreliable distributed infrastructures''. IEEE Trans.
  Parallel Distribut. Syst. vol 18, no 11 (2007)\\
\item
  Christoforou, E., Anta, A., Georgiou, C., Mosteiro, M.: Algorithmic
  mechanisms for reliable master-worker internet-based computing. IEEE
  Trans. Comput. 63(1), 179--195 (2014)\\
\item
  Estrada, T., Taufer, M., Reed, K.: ``Modeling Job Lifespan Delays in
  Volunteer Computing Projects,'' in 9th IEEE/ACM International
  Symposium on Cluster Computing and the Grid (2009)\\
\item
  Heien, E.M., Anderson, D.P., Hagihara, K.: Computing low latency
  batches with unreliable Workers in Volunteer
\end{enumerate}

Computing Environments. J. Grid Comput. 7(4), 501--518 (2009)\\
45. Lee, Y., Zomaya, A., Siegel, H.: ``Robust task scheduling for
volunteer computing systems''. J. Supercomput. vol 53, no 163 (2010)\\
46. Chernov, I.: ``Theoretical study of replication in desktop grid
computing: Minimizing the mean cost,'' in Proceedings of the 2nd
Applications in Information Technology (ICAIT-2016), Aizu-Wakamatsu,
Japan (2016)\\
47. Watanabe, K., Fukushi, M., Horiguchi, S.: ``Optimal Spot-checking to
minimize the Computation Time in Volunteer Computing,'' in Proceedings
of the 22nd IPDPS conference, PCGrid2008 workshop, Miami (2008)\\
48. Cano, P.P., Vargas-Lombardo, M.: Security threats in volunteer
computing environments using the Berkeley open infrastructure for
network computing (BOINC). Comput. Res. Model. 7(3), 727--734 (2015)\\
49. Elwaer, A., Harrison, A., Kelley, I., Taylor, I.: ``Attic: a case
study for distributing data in BOINC projects''. In IEEE Ninth
International Symposium on Parallel and Distributed Processing with
Applications Workshops, Busan, South Korea (2011)\\
50. Cole, N., Desell, T., González, D. L., Vega, F. F. D.,
Magdon-Ismail, M., Newberg, H., Szymanski, B., Varela, C.: ``Genetic
Algorithm Implementation for Boinc''. In de Vega F.F., Cantú-Paz E.
(eds) Parallel and Distributed Computational Intelligence. Studies in
Computational Intelligence, vol 269, pp 63--90 (2010)\\
51. Feki, M. S., Nguyen, V. H., Garbey, M.: ``Genetic algorithm
implementation for BOINC''. Adv. Parallel Comput. Vol 19, Elsevier, pp
212--219 (2010)\\
52. Nikitina, N., Ivashko, E., Tchernykh, A.: ``Congestion Game
Scheduling Implementation for High-Throughput Virtual Drug Screening
Using BOINC-Based Desktop Grid,'' in Malyshkin V. (eds) Parallel
Computing Technologies. PaCT 2017. Lecture Notes in Computer Science,
vol 10421., Nizhni Novgorod, Russia, Springer, pp 480--491 (2017)\\
53. Kanna, N., Subhlok, J., Gabriel, E., Rohit, E., Anderson, D.: ``A
Communication Framework for Fault-tolerant Parallel Execution,'' in The
22nd International Workshop on Languages and Compilers for Parallel
Computing, Newark, Delaware (2009)\\
54. Costa, F., Veiga, L., Ferreira, P.: ``BOINC-MR: MapReduce in a
Volunteer Environment,'' in Meersman R. et al.~(eds) On the Move to
Meaningful Internet Systems: OTM 2012, Springer (2012)\\
55. Schlitter, N., Laessig, J., Fischer, S., Mierswa, I., ``Distributed
data analytics using RapidMiner and BOINC,'' in Proceedings of the 4th
RapidMiner Community Meeting and Conference (RCOMM 2013), Porto,
Portugal (2013)\\
56. Desell, T.: ``Large scale evolution of convolutional neural networks
using volunteer computing,'' in Proceedings of the Genetic and
Evolutionary Computation Conference Companion, Berlin (2017)\\
57. Ivashko, E., Golovin, A.: ``Partition Algorithm for Association
Rules Mining in BOINC-based Enterprise Desktop Grid,'' in International
Conference on Parallel Computing Technologies, Petrozavodsk (2015)

\begin{enumerate}
\def\labelenumi{\arabic{enumi}.}
\setcounter{enumi}{57}
\item
  Vega-Rodríguez, M., Vega-Pérez, D., Gómez-Pulido, J., Sánchez-Pérez,
  J.: ``Radio Network Design Using Population-Based Incremental Learning
  and Grid Computing with BOINC,'' in Applications of Evolutionary
  Computing., Heidelber (2007)\\
\item
  Pellicer, S., Ahmed, N., Pan, Y., Zheng, Y.: ``Gene Sequence Alignment
  on a Public Computing Platform,'' in Proceedings of the 2005
  International Conference on Parallel Processing Workshops, Oslo
  (2005)\\
\item
  Cochran, E.S., Lawrence, J.F., Christensen, C., Jakka, R.S.: The
  quake-catcher network: citizen science expanding seismic horizons.
  Seismol. Res. Lett. 80(1), 26--30 (2009)\\
\item
  Desell, T., Bergman, R., Goehner, K., Marsh, R., VanderClute, R.,
  Ellis-Felege, S.: ``Wildlife@Home: Combining Crowd Sourcing and
  Volunteer Computing to Analyze Avian Nesting Video,'' in 2013 IEEE 9th
  International Conference on eScience, Beijing (2013)\\
\item
  Silva, J. N., Veiga, L., Ferreira, P.: ``nuBOINC: BOINC Extensions for
  Community Cycle Sharing,'' in Second IEEE International Conference on
  Self-Adaptive and Self-Organizing Systems Workshops, Venice, Italy
  (2008)\\
\item
  Binti, N. N., Zakaria, M. N. B., Aziz, I. B. A., Binti, N. S.:
  ``Optimizing BOINC scheduling using genetic algorithm based on thermal
  profile,'' in 2014 International Conference on Computer and
  Information Sciences, Kuala Lumpur, Malaysia (2014)\\
\item
  Science United, https://scienceunited.org/ (2018)
\item
  Klimeck, G., McLennan, M., Brophy, S. P., Adams, G. B., Lundstrom, M.
  S.: ``nanoHUB.org: Advancing Education and Research in
  Nanotechnology,'' Computing in Science \& Engineering, vol 10, no 17
  (2008)\\
\item
  Shepherd, C., Arfaoui, G., Gurulian, I., Lee, R., Markantonakis, K.,
  Akram, R., Sauveron, D., Conchon, E.: ``Secure and Trusted Execution:
  Past, Present, and Future - A Critical Review in the Context of the
  Internet of Things and Cyber-Physical Systems,'' in 2016 IEEE
  Trustcom/BigDataSE/ISPA (2016)\\
\item
  Nov, O., Arazy, O., Anderson, D.: ``Technology-Mediated Citizen
  Science Participation: A Motivational Model,'' in Fifth International
  AAAI Conference on Weblogs and Social Media (ICWSM 2011), Barcelona
  (2011)\\
\item
  Raoking, F., Cohoon, J. M., Cooke, K., Taufer, M., Estrada, T.:
  ``Gender and volunteer computing: A survey study,'' in IEEE Frontiers
  in Education Conference (FIE) Proceedings, Madrid (2014)\\
\item
  Gridcoin, ``The Computation Power of a Blockchain Driving Science and
  Data Analysis,'' https://gridcoin.us/assets/img/whitepaper.pdf
  (2018)\\
\item
  Chernov, I., Nikitina, N., Ivashka, E.: ``Task scheduling in desktop
  grids: open problems''. Open Eng. vol 7, no 1 (2017)\\
\item
  Dinis, G., Zakaria, N., Naono, K.: ``Pluggable scheduling on an
  open-source based volunteer computing infrastructure,'' in 2014
  International Conference on Computer and Information Sciences, Kuala
  Lumpur, Malaysia (2014)
\end{enumerate}

\end{document}
